\section{Local Optimization Method: MPPI}
\label{sec:local}

My local planner is a \textbf{Model Predictive Path Integral} controller. It is a sampling-based stochastic
optimal-control method: at every control step it evaluates many candidate motions against
a cost function and returns the best velocity command, continuously adapting to the nearby
obstacle geometry.

\paragraph{Motion model.} The robot state is $(r,c,\theta,v,\omega)$ and evolves under a
differential-drive model integrated at $\Delta t$:
\begin{equation}
    \theta \mathrel{+}= \omega\,\Delta t,\quad
    r \mathrel{+}= v\cos\theta\,\Delta t,\quad
    c \mathrel{+}= v\sin\theta\,\Delta t .
\end{equation}
The robot cannot translate sideways and its velocities are bounded, so these constraints
are respected directly in the rollout.

\paragraph{Sampling-based optimisation.} Rather than committing to a single constant
command, MPPI optimises a whole control sequence $U=(u_0,\dots,u_{T-1})$ over a short
horizon $T$. Each step it draws $K$ noisy rollouts $U^{(k)}=U+\mathcal{E}^{(k)}$ (clipped to
the velocity limits), simulates each with the motion model, scores each trajectory cost
$S(\tau^{(k)})$, and updates the nominal sequence by a softmin-weighted average of the
perturbations:
\begin{equation}
    w_k = \frac{\exp\!\big(-\tfrac1\lambda(S_k-\min_j S_j)\big)}{\sum_j \exp\!\big(-\tfrac1\lambda(S_j-\min_i S_i)\big)},
    \qquad U \leftarrow U + \sum_k w_k\,\mathcal{E}^{(k)} .
\end{equation}
It then executes the first command $u_0$ and shifts $U$ forward one step (receding horizon),
which warm-starts the next optimisation and yields smooth, temporally consistent commands.
Because it samples full sequences and softly averages over many rollouts, MPPI can represent
compound manoeuvres (``turn, then go straight'') and steer around local minima.

\paragraph{Objective function.} Each rollout is scored by a weighted sum (lower is better):
\begin{equation}
    S(\tau) = w_g\,c_\text{goal} + \alpha\,c_\text{head} + \beta\,c_\text{obs}
            + \gamma\,c_\text{vel} + \delta\,c_\text{smooth} + \epsilon\,c_\text{dev},
\end{equation}
whose terms reward closing distance to the goal, heading alignment, obstacle clearance, forward speed, and low angular jerk and staying
near the global path. The goal term is the primary driver of progress; heading alone only
aligns orientation and gives no incentive to actually advance.

\paragraph{Weighting, and speed/safety/efficiency trade-offs.} The weights encode the
robot's ``personality'' and expose direct trade-offs. A high $\beta$ (obstacle) yields wide,
cautious detours; a high $\gamma$ or $w_g$ yields fast but riskier passes close to obstacles;
a high $\epsilon$ keeps the robot glued to the global path at the cost of local shortcuts.
No single weighting is best --- the right balance of safety, speed and efficiency depends on
obstacle density --- and because MPPI re-solves this optimisation every control step against
the current geometry, it reacts to obstacles rather than blindly replaying the global path.
